\section{결론 및 향후 연구}
\label{sec:conclusion}

본 논문은 정간보의 격자 구조를 텍스트에 그대로 옮긴 경량 문법과, 이를
단일 원천으로 삼는 서버리스 웹 편집·조판 시스템 ``정간보 생성기''를
제안했다. 시스템은 WYSIWYG과 텍스트 편집의 양방향 동기화, 시김새 59종의
벡터 렌더링과 텍스트 직렬화, 전통 조판 관행의 규칙화, mm 단위 SVG 조판을
통한 인쇄·PDF·PNG 출력, JSON 문서 교환, Web Audio 재생을 정적 파일만으로
구현하여, 워드프로세서 수작업에 머물러 있던 정간보 제작을 데이터 기반
워크플로로 전환할 수 있음을 보였다.

향후 연구는 세 방향이다.
첫째, \textbf{소니피케이션 고도화}: 시김새의 음고 궤적과 장단 타점을
합성에 반영하고, 산조·정악 등 장르별 음색 모델을 얹어 ``들리는 정간보''로
확장한다.
둘째, \textbf{상호운용성}: \code{.jgb.json}과 MEI/MusicXML 간 변환기를
구현해 오선보 병기와 음악학 도구 연계를 지원하고, 정간보
OMR~\cite{kim2025omr} 출력의 표준 컨테이너로 문법을 제안한다.
이는 세종실록악보 등 고악보 아카이브~\cite{han2024sixdragons}를
편집 가능한 형태로 되살리는 파이프라인의 마지막 조각이 된다.
셋째, \textbf{사용자 연구}: 국악 교육 현장 투입을 통한 과제 기반 사용성
평가와 교육 효과 검증을 수행한다.

시스템은 오픈 웹에 공개되어 있다.
\TODO{배포 URL·저장소 공개 여부 결정 후 기입 (GitHub Pages 주소)}

% ------------------------------------------------------------
\section*{감사의 글}
\TODO{연구비 지원·사사 문구 (해당 시)}
